% Groupes caractéristiques entourés : éthanol et acide éthanoïque (corrigé de l'exercice 9, chapitre 8)
\begin{tikzpicture}[x=1cm,y=1cm,
  formule/.style={font=\fontsize{11}{13}\selectfont,inner sep=1pt,anchor=base west},
  groupe/.style={draw=otGreenDark,line width=0.8pt,rounded corners=1.5mm,dash pattern=on 2.4pt off 1.4pt,inner sep=2.5pt,
    font=\fontsize{11}{13}\selectfont,anchor=base west},
  nomG/.style={pcLegendeG,anchor=north},nomL/.style={pcLegende,anchor=north,text=otGreenDark}]
\node[formule] (a1) at (0,0) {$\ce{CH3-CH2-}$};
\node[groupe] (a2) at (a1.base east) {$\ce{OH}$};
\node[nomG] at (1.15,-0.45) {éthanol};
\node[nomL] at (1.15,-0.85) {groupe hydroxyle : alcool};
\draw[otGreenDark,line width=1pt,-{Stealth[length=2.4mm]}] (3.0,0.12) -- (4.4,0.12) node[midway,above,pcLegende,text=otGreenDark] {oxydation};
\node[formule] (b1) at (5.0,0) {$\ce{CH3-}$};
\node[groupe] (b2) at (b1.base east) {$\ce{COOH}$};
\node[nomG] at (6.0,-0.45) {acide éthanoïque};
\node[nomL] at (6.0,-0.85) {groupe carboxyle : acide carboxylique};
\end{tikzpicture}
\par\smallskip\legende{En s'oxydant, le carbone qui portait le groupe $\ce{-OH}$ devient celui du groupe carboxyle.}
